\documentclass[french]{llncs}

\usepackage[english, french]{babel}
\usepackage[utf8]{inputenc}

\usepackage{graphicx}

\usepackage{listings}
\usepackage{subfig}
\usepackage{hyperref}

%-----------------------------------------
\title{Titre de l'étude}

%\subtitle{\vspace{.5cm}Rapport de Master}

\author{ Prénom \textsc{Nom} \\
         Prénom \textsc{Nom} }

\institute{Université Claude Bernard Lyon 1, France}

\lstset{ 
    basicstyle=\small,
    captionpos=b  % sets the caption-position to bottom
}

\providecommand{\motcles}[1]{\textbf{Mots-clés:} #1}

\pagestyle{headings} 
%-----------------------------------------

\begin{document}

    \renewcommand{\thelstlisting}{\arabic{lstlisting}}

    \maketitle   
    
    \selectlanguage{english}
    \begin{abstract}
        Abstract in english.
        
         \begin{keywords}
        keywords
    \end{keywords}
    \end{abstract}


   

%-----------------------------------------

    \section{Introduction} \label{intro}

Pour produire votre rapport de projet en \LaTeX\, vous devez utiliser la
classe \textit{llncs.cls}. \\
        
        Toutes les instructions détaillées concernant l'utilisation de la classe \textit{llncs.cls}, ainsi que les templates \LaTeX. et \texttt{Word}, sont téléchargeables sur la page suivante:\\ \url{https://www.springer.com/gp/computer-science/lncs/conference-proceedings-guidelines}.\\
        
     Le rapport de projet doit faire maximum 4 pages (hors refs). Vous pouvez l'écrire en Français ou en Anglais mais l'abstract doit être en Anglais.
      
      
    \section{Section de niveau 1} \label{sect:niveau11}

Exemple de section de niveau 1.  Voici un premier paragraphe. Voici
un premier paragraphe, Voici un premier paragraphe, Voici un premier
paragraphe, Voici un premier paragraphe, Voici un premier
paragraphe, Voici un premier paragraphe, Voici un premier
paragraphe, Voici un premier paragraphe, Voici un premier
paragraphe.

       

        \subsection{Section de niveau 2} \label{sect:niveau2}

           Exemple de section de niveau 2.

      \subsubsection{Section de niveau 3} \label{sect:niveau3}

  Exemple de section de niveau 3.\\
  
 \noindent Vous pouvez aussi utiliser 
\verb+\paragraph{+\textit{Titre du paragraphe.}\verb+}+ qui produit

\paragraph{Titre du paragraphe.} Ceci est un exemple de section type
paragraphe. On note l'absence de numérotation.
  
       \subsection{Liste d'éléments}
         
      Exemple de liste sans numérotation :
\begin{itemize}
    \item premier item ;
    \item deuxième item ;
    \begin{itemize}
      \item premier sous-item,
      \item second sous-item,
      \item dernier sous-item.
  \end{itemize}
  \item dernier item.
\end{itemize}

Autre liste possible avec numérotation :
\begin{enumerate}
    \item premier item ;
    \item deuxième item ;
    \begin{enumerate}
      \item premier sous-item,
      \item second sous-item,
      \item dernier sous-item.
    \end{enumerate}
  \item dernier item.
\end{enumerate}

\subsection{Tables, Figures}


Les tableaux et figures doivent être centrés horizontalement. Les
légendes sont placées sous ces mêmes éléments et sont centrées si
elle font moins d'une ligne et justifiées sinon. Voir le
tableau~\ref{tab_exemple} et la figure~\ref{fig_exemple} pour des
exemples. Les légendes doivent se terminer avec un point.

\begin{table}[ht]
 \begin{center}
   \begin{tabular}{l | c c }
   \hline
                & modèle 1 & modèle 2 \\
   \hline
   indicateur 1 & 1        & 2        \\
   indicateur 2 & 2        & 1        \\
   indicateur 3 & 1        & 1        \\
   \hline
   \end{tabular}
\caption{Exemple de tableau avec sa légende.} \label{tab_exemple}
 \end{center}
\end{table}


\begin{figure}[t]
\begin{center}
 	\includegraphics[width=0.5\linewidth]{logoucbl.png}
 \caption{Exemple de figure avec sa légende. La légende est centrée,
sauf si, comme ici, elle fait plus d'une ligne, auquel cas elle est
justifiée.} \label{fig_exemple}
\end{center}
\end{figure}

   

       \subsection{Equation} \label{sect2:eq}

Pour inclure une équation, vous pouvez utiliser l'environnement \texttt{equation} commue pour l'exemple ci-dessous: 
\begin{equation}
\left(\frac{a^{2} + b^{2}}{c^{3}} \right) = 1 \quad
\mbox{ if } c\neq 0 \mbox{ and if } a,b,c\in \bbbr \enspace .
\end{equation}
           
           \subsection{Notes de bas de pages}

Le texte de la note apparaît en bas de la page.\footnote{Exemple de
note de bas de page.}


    \section{Consignes pour les références} \label{sect3}

Les références sont données en fin de rapport.

Nous vous recommandons l'utilisation de
\textsc{Bib}\TeX\ avec le style bibliographique \textit{splncs03.bst}. Par exemple pour ce document, nous utilisons le fichier \texttt{bibM2.bib}, dans lequel est présent la référence \cite{Russell2003}. 
     
Dans le corps du texte, on utilise la commande \texttt{$\backslash$ cite} pour faire référence à un article.  

    \newpage
    \bibliographystyle{splncs03} 
    \bibliography{bibM2.bib}        

\end{document}
